\section{Synthesis: what no family governs as a whole}
\label{sec:synthesis}

\Cref{tab:matrix} summarises where each family decides and enforces across the
stages of a multi-source study. A cell names the usual practice in the family,
not every system in it.

\begin{table}[t]
\centering
\footnotesize
\setlength{\tabcolsep}{3pt}
\begin{tabularx}{\linewidth}{@{}p{0.14\linewidth} Y Y Y Y Y p{0.13\linewidth}@{}}
\toprule
Family & Decision made by & Read & Train or compute & Send a result & Release a model & Planner fails \\
\midrule
F1 Responsible-AI practice & people; classifiers for LLM text & LLM input and output only & design time, by people & -- & audit and documentation, by people & often open (Guardrails AI: \code{noop}) \\
F2 Policy engines & policy language or engine; some analysable & permit or deny; transformations inside one warehouse & IFC in a runtime & egress checks in IFC systems & -- & varies by product \\
F3 Privacy techniques & a person sets \(\varepsilon\); software derives the mechanism & mechanism at output, central & DP training & local DP on devices & budget management & budget spent: refuse \\
F4 Federated learning & code, site files, people, ledgers, XACML & site filters (FLARE) & mechanisms in the protocol & site filters; secure aggregation & organisational, by people & rarely stated (FLA3: closed) \\
F5 Code to data & function libraries, checkers, vendor engines, compilers & function library at each site & site review of code & thresholds; output checkers & separate, later (GRAIMatter) & some closed (DataSHIELD) \\
F6 Architecture & planners and declarative specifications & mediators; sidecars & -- & -- & deployment, not models & varies (Envoy closed, Gatekeeper open) \\
F7 Release and regulation & pipeline steps, people, law & law states principles & -- & -- & gates in pipelines; people approve & not stated \\
\midrule
GRAILS & a planner over three inputs & per item, one vault & -- & -- & -- & not stated \\
Dagents & a total, typed planner & per field, at each site & \code{before\_train}; guarded payload & \code{before\_send}: clip or noise & gates; release Guard; a person & closed, recorded \\
\bottomrule
\end{tabularx}
\caption{Where each family decides and enforces across the stages of a
multi-source study. ``--'' means the family does not usually address the stage.}
\label{tab:matrix}
\end{table}

Five gaps follow. None of them means that no system does the thing; each
subsection of \cref{sec:survey} names systems that do part of it. The gaps are
in the combination.

\begin{description}[font=\normalfont\bfseries, leftmargin=1.2em, style=nextline]
  \item[G1. The form of the data, decided outside the store]
        Most decision points answer permit or deny. Where a system decides how
        much of the data to return and in what form, it usually decides inside
        the engine that already holds the pooled data: a warehouse, a clean room
        or a trusted research environment (F2, F3, F5).
  \item[G2. Enforcement at independent sources]
        Systems that compute at the source control what leaves through function
        libraries, people or vendor engines (F4, F5). Few apply one decision
        model at each source under a coordinator that may not read the source's
        rows.
  \item[G3. One rule set across the lifecycle]
        Reading, training, sending and releasing are governed by different
        mechanisms and often different people. Model release is the stage
        covered least, and it is usually governed by organisational rules
        (F4, F5, F7). DataSHIELD and Confidential Federated Computations are the
        exceptions we found.
  \item[G4. Decisions that can be read, tested and replayed]
        The rule about what leaves is often spread across server options,
        function code, checker guidance and vendor defaults (F5). Cedar's
        authorisation decisions and Conclave's and SMCQL's query plans are
        analysable (F2, F5), but they do not cover the form of the data, model
        updates or release together.
  \item[G5. Failure behaviour and evidence]
        What happens when the decision point is unreachable is rarely
        documented, and defaults differ: some products deny, others admit
        (F1, F2, F6). Logs record decisions, but they rarely name the logic that
        decided, and no regulation asks for that per record, although the AI
        Act's record-keeping duty supports it (F7).
\end{description}

GRAILS addresses G1 for one application with one store. The design described in
the next section is one answer to G1 to G5 together. It is built from known
primitives: the PDP/PEP split, purpose binding, minimum-frequency rules,
clipping, release gates and hash chains. What it adds is their combination in
one typed decision function, enforced at four boundaries at each source.
